\chapter{Dữ liệu và tiền xử lý dữ liệu}

\section{Nguồn dữ liệu}

Dữ liệu của hệ thống gồm hai nhóm chính: dữ liệu do người dùng cung cấp và dữ liệu do hệ thống tự thu thập hoặc tự suy diễn.

% \begin{table}[H]
%   \centering
%   \caption{Các nhóm dữ liệu sử dụng trong hệ thống}
%   \begin{tabular}{p{4cm} p{9cm}}
%     \toprule
%     \textbf{Nhóm dữ liệu} & \textbf{Mô tả} \\
%     \midrule
%     Thông tin hộ gia đình & Số điện tháng trước, số người trong hộ, vị trí hoặc tỉnh/thành phố. \\
%     Thiết bị điện & Danh sách thiết bị, số lượng từng thiết bị, nhóm thiết bị và công suất ước lượng. \\
%     Thời gian & Tháng dự báo, số ngày trong tháng, số ngày cuối tuần, mùa trong năm. \\
%     Khí tượng & Nhiệt độ trung bình, độ ẩm trung bình, số ngày nóng, số ngày mưa. \\
%     Nhãn đầu ra & Điện năng tiêu thụ thực tế của tháng cần dự báo, đơn vị kWh. \\
%     \bottomrule
%   \end{tabular}
% \end{table}

\section{Thông tin người dùng cung cấp}

Người dùng cần cung cấp tối thiểu các thông tin sau:

\begin{itemize}[leftmargin=1.2cm]
  \item \texttt{kWh\_previous}: số điện tiêu thụ của tháng trước.
  \item \texttt{num\_people}: số người sinh sống trong hộ.
  \item Danh sách thiết bị điện trong nhà thông qua giao diện chọn thiết bị.
  % \item Vị trí hoặc tỉnh/thành phố để hệ thống lấy dữ liệu thời tiết.
\end{itemize}

\section{Thiết kế dữ liệu thiết bị}

Thay vì chỉ sử dụng tổng số thiết bị, hệ thống chuyển danh sách thiết bị thành các đặc trưng tổng hợp có ý nghĩa hơn. Mỗi thiết bị có thể được mô tả bằng các thuộc tính:

\begin{table}[H]
  \centering
  \caption{Cấu trúc thông tin cho một thiết bị điện}
  \begin{tabular}{p{4cm} p{9cm}}
    \toprule
    \textbf{Thuộc tính} & \textbf{Ý nghĩa} \\
    \midrule
    \texttt{device\_name} & Tên thiết bị, ví dụ điều hòa, tủ lạnh, máy giặt. \\
    \texttt{category} & Nhóm thiết bị như làm mát, nhà bếp, giải trí, tải nền. \\
    \texttt{quantity} & Số lượng thiết bị người dùng chọn. \\
    \texttt{default\_power\_w} & Công suất trung bình ước lượng theo watt. \\
    \texttt{usage\_level} & Mức sử dụng mặc định: thấp, trung bình, cao hoặc luôn bật. \\
    \bottomrule
  \end{tabular}
\end{table}

Từ dữ liệu thiết bị, hệ thống sinh ra các đặc trưng:

\begin{itemize}[leftmargin=1.2cm]
  \item \texttt{num\_devices}: tổng số thiết bị.
  \item \texttt{total\_device\_power}: tổng công suất ước lượng.
  \item \texttt{high\_power\_device\_count}: số thiết bị công suất lớn.
  \item \texttt{cooling\_device\_count}: số thiết bị làm mát.
  \item \texttt{always\_on\_power}: công suất của nhóm thiết bị thường xuyên bật.
  \item \texttt{estimated\_device\_kWh}: điện năng ước lượng từ công suất và thời gian sử dụng mặc định.
\end{itemize}

\section{Dữ liệu thời gian, vị trí và khí tượng}

Sau khi người dùng cho phép truy cập vị trí, hệ thống có thể gọi API thời tiết hoặc sử dụng dữ liệu lịch sử để lấy các đặc trưng khí tượng. Các đặc trưng đề xuất gồm:

\begin{itemize}[leftmargin=1.2cm]
  \item \texttt{month}: tháng cần dự báo.
  \item \texttt{season}: mùa trong năm hoặc nhóm mùa theo khí hậu Việt Nam.
  \item \texttt{avg\_temperature}: nhiệt độ trung bình tháng.
  \item \texttt{avg\_humidity}: độ ẩm trung bình tháng.
  \item \texttt{hot\_days}: số ngày nhiệt độ vượt ngưỡng nóng.
  \item \texttt{rainy\_days}: số ngày mưa.
\end{itemize}

\section{Bộ dữ liệu huấn luyện đề xuất}

Mỗi dòng dữ liệu huấn luyện tương ứng với một hộ gia đình tại một tháng cụ thể. Cấu trúc dữ liệu đề xuất:

\begin{longtable}{p{5cm} p{8cm}}
  \caption{Các thuộc tính trong bộ dữ liệu huấn luyện đề xuất} \\
  \toprule
  \textbf{Thuộc tính} & \textbf{Ý nghĩa} \\
  \midrule
  \endfirsthead
  \toprule
  \textbf{Thuộc tính} & \textbf{Ý nghĩa} \\
  \midrule
  \endhead
  \texttt{household\_id} & Mã hộ gia đình. \\
  \texttt{month} & Tháng quan sát. \\
  \texttt{province} & Tỉnh hoặc thành phố. \\
  \texttt{kWh\_previous} & Điện năng tiêu thụ tháng trước. \\
  \texttt{num\_people} & Số người trong hộ. \\
  \texttt{num\_devices} & Tổng số thiết bị. \\
  \texttt{total\_device\_power} & Tổng công suất thiết bị ước lượng. \\
  \texttt{cooling\_device\_count} & Số thiết bị làm mát. \\
  \texttt{avg\_temperature} & Nhiệt độ trung bình tháng. \\
  \texttt{avg\_humidity} & Độ ẩm trung bình tháng. \\
  \texttt{hot\_days} & Số ngày nắng nóng trong tháng. \\
  \texttt{kWh\_next} & Điện năng tiêu thụ thực tế của tháng tiếp theo. \\
  \bottomrule
\end{longtable}

\section{Tiền xử lý dữ liệu}

% Quá trình tiền xử lý gồm các bước:

% \begin{enumerate}[leftmargin=1.2cm]
%   \item Kiểm tra và loại bỏ dữ liệu thiếu hoặc bất thường.
%   \item Chuẩn hóa đơn vị đo, ví dụ điện năng theo kWh, công suất theo W.
%   \item Mã hóa các biến phân loại như tháng, mùa, tỉnh/thành phố.
%   \item Chuẩn hóa các biến liên tục bằng Min-Max Scaling hoặc Standard Scaling.
%   \item Chia dữ liệu thành tập huấn luyện, tập kiểm định và tập kiểm tra.
% \end{enumerate}

\section{Chuẩn hóa dữ liệu}

Với ANFIS, việc chuẩn hóa dữ liệu giúp quá trình học ổn định hơn, đặc biệt khi các biến có miền giá trị rất khác nhau. Công thức Min-Max Scaling:

\begin{equation}
  x' = \frac{x - x_{\min}}{x_{\max} - x_{\min}}
\end{equation}

Sau khi dự báo, kết quả có thể được chuyển ngược về đơn vị kWh để hiển thị cho người dùng.
